\documentclass{article}
\usepackage{graphicx} % Required for inserting images
\usepackage{amsmath}
\usepackage{bm}

\title{Falsifier-Based robustification for OpenSCVX}
\date{October 2026}

\begin{document}

\maketitle

\section{Introduction}
Robust trajectory optimization is vital for autonomous robotics systems operating under uncertain dynamics and environmental disturbances. While successive convexification (SCVx) enables efficient nonlinear trajectory optimization, nominal solutions may remain vulnerable to disturbances. The hypothesis is we propose a counterexample-guided robustification framework for OpenSCvx that integrates adversarial falsification with trajectory optimization under one system. A falsifier identifies disturbances causing continuous-time safety violations, and those disturbances are explicitly converted into constraints for subsequent SCvx iterations. What we must prove are two matters: (1) demonstrate a competitive robustness-efficiency tradeoff against existing optimization methods (2) establish whether explicit constraint-based robustification offers practical advantages over neural policy fine-tuning.
\section{State-of-the-Art}
\begingroup
\setlength{\parskip}{0.8em}
\setlength{\parindent}{0pt}

Mutapcic and Boyd~\cite{mutapcic2009} developed an optimization method that repeated searches for worst-case scenarios and add new constraints to prevent similar failures. Dawson and Fan~\cite{dawson2022} applied a related approach to robotic trajectory planning where an optimizer learns from failures identified by a falsifier. Their method produced more robust trajectories while taking less than half the time of competing methods in their experiments.

Shi et al.~\cite{shi2024} used RL (PPO) to discover realistic disturbances that could cause quadrupedal robots to fail. They then retrained the robot's neutral controller's using these failures to improve its robustness. 

Abdul et el.~\cite{abdul2025} developed a SCvx method that plans trajectories while accounting for all disturbances within a predefined bounded set. Their approach also accounts for errors caused by approximating nonlinear dynamics. Wang et al.~\cite{wang2026} later accelerated this method using GPUs, achieving speedups of up to 139.6 times over the original implementation. Elango et al. ~\cite{elango2025} developed a way for SCvx to enforce safety constraints throughout the entire trajectory rather than only at selected time steps. 

Building on to these works, we adding a falsifier to OpenSCvx that actively searches
for disturbances that cause trajectory failures. Each discovered failure introduces a new safety constraint and allow SCvx to replan and avoid similar failures. We aim to investigate whether this approach can achieve comparable robustness with lower computational cost and make safety improvements easier to understand. 



\section{Limitations}
Existing robust SCvx methods~\cite{abdul2025,wang2026} account for bounded disturbances, but solving these optimization problems can be computationally heavy. Counterexample-guided methods~\cite{dawson2022} improve planning by learning from failures, but do not directly integrate hard continuous-time safety constraints into SCvx. Meanwhile, adversarial RL~\cite{shi2024} improves neutral controllers through fine-tuning rather than explicit constraint updates. Potential benefits of combining adaptive falsification with continuous-time SCvx remain unexplored in the works reviewed here. 
\section{Ideas}
A unicycle model corresponds to a rover navigating a surface environment, and this has a wide range of applications, including, but not limited to, the following case studies: 
\begin{itemize}
\item Surface exploration of Moon/Mars (the model corresponds to the physical system)
\item Surface exploration of unknown Earth environments (the model simplifies a legged robot, but may be representative for planning)
\item Planning of autonomous vehicles in traffic: e.g. car, trucks... (the model simplifies a car, but may be representative for planning)
\end{itemize}
\section{Mathematical Mature Formulation}

We consider a planar unicycle with continuously actuated
forward and angular velocities. The robot moves without
lateral slip, and its dynamics are deterministic.
The planning horizon $\mathcal H=[0,T]$ is fixed.
Initial and terminal configurations are prescribed.

The state and control vectors are
\begin{equation}
\bm{x} =
\begin{bmatrix}p_x\\p_y\\\theta\end{bmatrix},
\qquad
\bm{u} =
\begin{bmatrix}v\\\omega\end{bmatrix}.
\end{equation}

The nonlinear dynamics are
\begin{equation}
\bm{f}(\bm{x},\bm{u}) =
\begin{bmatrix}
v\cos\theta\\
v\sin\theta\\
\omega
\end{bmatrix}.
\end{equation}

Let $\bm{g}(\bm{x})\leq\bm{0}$ denote the
collision-avoidance constraints and
$\bm{G}(\bm{u})\leq\bm{0}$ the actuator limits.
The objective is to minimize the traveled distance.

\medskip
\noindent\textbf{Problem 1 (P1)}

\begin{equation}
\begin{aligned}
\underset{\bm{x}(\cdot),\,\bm{u}(\cdot)}
{\operatorname{minimize}}
\quad & \int_0^T v(t)\,dt
\\[4pt]
\text{subject to}\quad &
\left|\;\;
\begin{aligned}
\dot{\bm{x}} &= \bm{f}(\bm{x},\bm{u})\\
\bm{g}(\bm{x}) &\leq \bm{0}
\\
\bm{G}(\bm{u}) &\leq\bm{0}
\\
\bm{x}(0) &= \bm{x}_{\mathrm{init}}
\\
\bm{x}(T) &= \bm{x}_{\mathrm{goal}}
\end{aligned}
\right.
\end{aligned}
\label{eq:unicycle_ocp}
\end{equation}

The actuator limits are defined by
\begin{equation}
\bm{G}(\bm{u}) =
\begin{bmatrix}
-v\\
v-v_{\max}\\
\omega-\omega_{\max}\\
-\omega-\omega_{\max}
\end{bmatrix}.
\end{equation}


\subsection{Counterexample-Guided Robustification}

We extend Problem 1 with a bounded constant steering
disturbance $b\in\mathcal B=[-\bar b,\bar b]$:
\begin{equation}
\dot{\bm{x}}_b =
\begin{bmatrix}
v\cos\theta_b\\
v\sin\theta_b\\
\omega+b
\end{bmatrix}.
\end{equation}

\paragraph{Adversarial Falsification.}
For a circular obstacle with center $\bm c$ and
radius $R$, define the violation function
\begin{equation}
q(\bm{x}_b)=R^2-\|\bm p_b-\bm c\|_2^2.
\end{equation}

Given the current control $\bm u^j$, the falsifier
searches for the worst-case disturbance and time:
\begin{equation}
\boxed{
(b_j^*,t_j^*)\in
\underset{b\in\mathcal B,\ t\in[0,T]}
{\operatorname{argmax}}
q(\bm{x}_b(t;\bm u^j)).
}
\end{equation}

A positive maximum indicates a safety violation.

\paragraph{Counterexample Feedback.}
Starting with $\mathcal S_0=\{0\}$, each discovered
violating disturbance is retained:
\begin{equation}
\boxed{
\mathcal S_{j+1}=\mathcal S_j\cup\{b_j^*\}.
}
\end{equation}

The next optimization introduces one state trajectory
per retained disturbance, with shared controls:
\begin{equation}
\begin{aligned}
\min_{\bm u,\{\bm{x}_b\}}\quad
&\int_0^T v(t)\,dt\\
\mathrm{s.t.}\quad
&\dot{\bm{x}}_b=\bm f_b(\bm{x}_b,\bm u),
&&b\in\mathcal S_{j+1},\\
&q(\bm{x}_b(t))\leq0,
&&\forall t,\ b\in\mathcal S_{j+1},\\
&\bm G(\bm u(t))\leq\bm0,
&&\forall t\in[0,T],\\
&\bm{x}_b(0)=\bm{x}_{\mathrm{init}},
&&b\in\mathcal S_{j+1},\\
&\bm{x}_0(T)=\bm{x}_{\mathrm{goal}}.
\end{aligned}
\end{equation}

\paragraph{Integration into SCvx.}
For each scenario, SCvx linearizes the disturbed
dynamics around a reference trajectory and applies
the affine obstacle-avoidance restriction:
\begin{equation}
\boxed{
q(\bar{\bm p}_b)
-2(\bar{\bm p}_b-\bm c)^\top
(\bm p_b-\bar{\bm p}_b)\leq0.
}
\end{equation}

This is a conservative convex restriction since
$q(\bm p_b)$ equals the affine expression minus
$\|\bm p_b-\bar{\bm p}_b\|_2^2$.

Continuous-time safety constraints can be handled
using OpenSCvx's existing framework.

The outer loop alternates falsification and SCvx
re-optimization, adding scenarios only when a
violation is discovered.

\paragraph{Scope.}
Only nominal terminal-state feasibility is imposed.
Safety over the full uncertainty set requires
certified falsification or additional verification.
Convergence and computational advantages of the
combined algorithm remain to be established.

\section{Development and Problem Solution}
OpenSCvx solves nonlinear optimal control problems by iteratively writing and
solving convex subproblems around a nominal trajectory. 

\begin{enumerate}
    \item A planar unicycle moves over a fixed \(8\,\mathrm{s}\) horizon.
    \item The vehicle starts at \((0,0,0)\) and must reach \((4,0,0)\).
    \item The controls satisfy \(0\le v\le1\,\mathrm{m/s}\) and \(|\omega|\le0.8\,\mathrm{rad/s}\).
    \item A circular obstacle centered at \((2,0)\) has radius \(0.55\,\mathrm{m}\).
    \item The path must satisfy \(\|\bm{p}(t)-(2,0)^\top\|_2\ge0.55\,\mathrm{m}\) for all \(t\in[0,8]\).
\end{enumerate}
A plot of the OpenSCvx solution follows.

\begin{figure}[h!]
\centering
\includegraphics[width=\textwidth]{unicycle_solution.png}
\caption{Solution to the unicycle problem}
\end{figure}

\newpage
\section{Conclusions}

\section{References}

% Add \usepackage{hyperref} in your preamble.

\begin{thebibliography}{9}

\bibitem{mutapcic2009}
A. Mutapcic and S. Boyd,
``Cutting-set methods for robust convex optimization
with pessimizing oracles,''
\textit{Optimization Methods and Software},
vol. 24, no. 3, pp. 381--406, 2009.
\url{https://doi.org/10.1080/10556780802712889}

\bibitem{dawson2022}
C. Dawson and C. Fan,
``Robust Counterexample-guided Optimization
for Planning from Differentiable Temporal Logic,''
arXiv:2203.02038, 2022.
\url{https://arxiv.org/abs/2203.02038}

\bibitem{shi2024}
F. Shi et al.,
``Rethinking Robustness Assessment:
Adversarial Attacks on Learning-based
Quadrupedal Locomotion Controllers,''
\textit{Robotics: Science and Systems}, 2024.
\url{https://doi.org/10.15607/RSS.2024.XX.057}

\bibitem{abdul2025}
A. T. Abdul, A. D. Saravanos,
and E. A. Theodorou,
``Nonlinear Robust Optimization for
Planning and Control,''
\textit{IEEE Conference on Decision
and Control (CDC)}, 2025.
\url{https://doi.org/10.1109/CDC57313.2025.11312909}

\bibitem{elango2025}
P. Elango et al.,
``Continuous-time successive convexification
for constrained trajectory optimization,''
\textit{Automatica},
vol. 180, article 112464, 2025.
\url{https://doi.org/10.1016/j.automatica.2025.112464}

\bibitem{wang2026}
J. Wang, A. T. Abdul, and E. A. Theodorou,
``cuNRTO: GPU-Accelerated Nonlinear Robust
Trajectory Optimization,''
\textit{Robotics: Science and Systems}, 2026.
\url{https://arxiv.org/abs/2603.02642}

\end{thebibliography}


\end{document}
